\documentclass[12pt,oneside]{book}

%============================================================
% PACKAGES
%============================================================

%--------------------
% Page + font
%--------------------
\usepackage[margin=1in]{geometry}
\usepackage[T1]{fontenc}
\usepackage[utf8]{inputenc}
\usepackage{lmodern}
\usepackage{microtype}
\usepackage{parskip}


%--------------------
% Math
%--------------------
\usepackage{amsmath, amssymb, amsfonts, amsthm, mathtools}
\usepackage{bm}
\usepackage{bbm}
\usepackage{mathrsfs}
\usepackage{dsfont}

%--------------------
% Tables + figures
%--------------------
\usepackage{graphicx}
\usepackage{float}
\usepackage{booktabs}
\usepackage{array}
\usepackage{xltabular}
\usepackage{multirow}
\usepackage{caption}
\usepackage{subcaption}

%--------------------
% Lists
%--------------------
\usepackage{enumitem}
\setlist[itemize]{topsep=2pt,itemsep=2pt}
\setlist[enumerate]{topsep=2pt,itemsep=2pt}

%--------------------
% Colors + hyperlinks
%--------------------
\usepackage[dvipsnames]{xcolor}
\usepackage[
    colorlinks=true,
    linkcolor=MidnightBlue,
    citecolor=MidnightBlue,
    urlcolor=MidnightBlue
]{hyperref}
\usepackage[nameinlink,capitalise]{cleveref}

%--------------------
% Algorithms
%--------------------
\usepackage{algorithm}
\usepackage{algpseudocode}

%--------------------
% Code blocks
%--------------------
\usepackage{listings}

\lstdefinestyle{codestyle}{
    basicstyle=\ttfamily\small,
    breaklines=true,
    frame=single,
    numbers=left,
    numberstyle=\tiny,
    keywordstyle=\color{MidnightBlue},
    commentstyle=\color{ForestGreen},
    stringstyle=\color{BrickRed},
    showstringspaces=false,
    tabsize=4
}

\lstset{style=codestyle}

%--------------------
% Section styling
%--------------------
\usepackage{titlesec}
\usepackage{tocloft}

\setcounter{tocdepth}{2}
\setcounter{secnumdepth}{3}

%============================================================
% THEOREMS
%============================================================

\numberwithin{equation}{section}

\theoremstyle{definition}
\newtheorem{definition}{Definition}[section]
\newtheorem{example}{Example}[section]
\newtheorem{exercise}{Exercise}[section]
\newtheorem{remark}{Remark}[section]

\theoremstyle{plain}
\newtheorem{theorem}{Theorem}[section]
\newtheorem{lemma}{Lemma}[section]
\newtheorem{proposition}{Proposition}[section]
\newtheorem{corollary}{Corollary}[section]

\theoremstyle{remark}
\newtheorem*{note}{Note}

%============================================================
% CUSTOM COMMANDS
%============================================================

% Sets
\newcommand{\R}{\mathbb{R}}
\newcommand{\N}{\mathbb{N}}
\newcommand{\Z}{\mathbb{Z}}
\newcommand{\Q}{\mathbb{Q}}
\newcommand{\C}{\mathbb{C}}

% Probability
\newcommand{\E}{\mathbb{E}}
\newcommand{\Var}{\mathrm{Var}}
\newcommand{\Cov}{\mathrm{Cov}}
\newcommand{\Corr}{\mathrm{Corr}}
\newcommand{\Prob}{\mathbb{P}}

\newcommand{\ind}{\mathbbm{1}}
\newcommand{\iid}{\overset{\text{iid}}{\sim}}
\newcommand{\convd}{\overset{d}{\to}}
\newcommand{\convp}{\overset{p}{\to}}
\newcommand{\as}{\overset{a.s.}{\to}}

% Operators
\newcommand{\given}{\,\middle|\,}
\newcommand{\argmin}{\operatorname*{arg\,min}}
\newcommand{\argmax}{\operatorname*{arg\,max}}

% Linear algebra / analysis
\newcommand{\norm}[1]{\left\lVert #1 \right\rVert}
\newcommand{\abs}[1]{\left\lvert #1 \right\rvert}
\newcommand{\inner}[2]{\left\langle #1, #2 \right\rangle}

\newcommand{\dd}{\,\mathrm{d}}
\newcommand{\grad}{\nabla}
\newcommand{\Hess}{\nabla^2}

\newcommand{\tr}{\operatorname{tr}}
\newcommand{\rank}{\operatorname{rank}}
\newcommand{\diag}{\operatorname{diag}}

% Distributions
\newcommand{\Normal}{\mathcal{N}}
\newcommand{\Bern}{\mathrm{Bernoulli}}
\newcommand{\Bin}{\mathrm{Binomial}}
\newcommand{\Pois}{\mathrm{Poisson}}
\newcommand{\Gam}{\mathrm{Gamma}}
\newcommand{\BetaDist}{\mathrm{Beta}}
\newcommand{\Exp}{\mathrm{Exponential}}
\newcommand{\Dir}{\mathrm{Dirichlet}}
\newcommand{\InvGam}{\mathrm{Inv\text{-}Gamma}}

% Title block
\title{EN.553.761 Unconstrained Nonlinear Optimization}
\author{Jonathan Y. Ma}
\date{June 2026}

\begin{document}

\maketitle

\noindent \emph{This course was taught by Dr. Shiqian Ma in the Fall 2026 Semester. The notes are transcribed by Jonathan Ma and may contain errors and omitted content, so any issues are to be directed towards me.}

\tableofcontents

\section{Gradient Method}

Consider
\[
    \min_{x\in\R^d} f(x),
\]
where $f$ is differentiable. A point $x^*$ is a global minimizer if
$f(x^*)\leq f(x)$ for every $x$, a local minimizer if this holds in a
neighborhood, and a stationary point if $\grad f(x^*)=0$.

\begin{definition}[Smoothness and strong convexity]
The function $f$ is $L$-smooth if
\[
    \norm{\grad f(x)-\grad f(y)}\leq L\norm{x-y}.
\]
It is $\mu$-strongly convex if
\[
 f(y)\geq f(x)+\inner{\grad f(x)}{y-x}
       +\frac{\mu}{2}\norm{y-x}^2.
\]
\end{definition}

For twice differentiable $f$, the sufficient Hessian conditions are
$\Hess f(x)\preceq LI$ for smoothness and
$\Hess f(x)\succeq\mu I$ for strong convexity. Convexity makes every local
minimum global; strong convexity makes the minimizer unique.

\begin{lemma}[Descent lemma]
If $f$ is $L$-smooth, then
\[
 f(y)\leq f(x)+\inner{\grad f(x)}{y-x}
            +\frac L2\norm{y-x}^2.
\]
\end{lemma}

\begin{proof}
Apply the fundamental theorem of calculus to
$f(x+t(y-x))$ and bound the change in the gradient using Lipschitz
continuity.
\end{proof}

Gradient descent uses
\begin{equation}
    x_{k+1}=x_k-\alpha_k\grad f(x_k).
    \label{eq:gradient-descent}
\end{equation}
With $\alpha=1/L$, the descent lemma gives
\[
    f(x_{k+1})\leq f(x_k)-\frac1{2L}\norm{\grad f(x_k)}^2.
\]

\begin{theorem}[Basic gradient-descent rates]
Let $f$ be convex and $L$-smooth, and let $x^*$ be a minimizer. Gradient
descent with $\alpha=1/L$ satisfies
\[
 f(x_k)-f(x^*)\leq\frac{L\norm{x_0-x^*}^2}{2k}.
\]
If $f$ is also $\mu$-strongly convex, then
\[
 f(x_k)-f(x^*)\leq
 \left(1-\frac{\mu}{L}\right)^k\{f(x_0)-f(x^*)\}.
\]
\end{theorem}

The condition number $\kappa=L/\mu$ controls the linear rate. Poor feature
scaling produces elongated level sets and slow zig-zagging. Stop using a
combination of gradient norm, relative objective change, step norm, and an
iteration budget; objective change alone can be misleading on flat regions.

\section{Quasi-Newton Method}

Newton's method locally approximates
\[
 f(x+p)\approx f(x)+g^\mathsf Tp+\tfrac12p^\mathsf THp
\]
and solves $Hp=-g$, where $g=\grad f(x)$ and $H=\Hess f(x)$. Near a
nondegenerate minimizer, full Newton steps converge quadratically under
standard Lipschitz-Hessian assumptions. Far from a minimizer, $H$ can be
indefinite and the Newton direction need not descend.

Quasi-Newton methods replace the Hessian or inverse Hessian with an estimate.
Set $s_k=x_{k+1}-x_k$ and $y_k=\grad f(x_{k+1})-\grad f(x_k)$. The secant
condition is
\[
    B_{k+1}s_k=y_k.
\]
The BFGS update is
\[
 B_{k+1}=B_k-\frac{B_ks_ks_k^\mathsf TB_k}{s_k^\mathsf TB_ks_k}
              +\frac{y_ky_k^\mathsf T}{y_k^\mathsf Ts_k}.
\]
If $B_k\succ0$ and $y_k^\mathsf Ts_k>0$, then $B_{k+1}\succ0$. Wolfe line
search guarantees this curvature condition for a descent direction.

In inverse form, with $\rho_k=(y_k^\mathsf Ts_k)^{-1}$,
\[
 H_{k+1}=(I-\rho_ks_ky_k^\mathsf T)H_k
         (I-\rho_ky_ks_k^\mathsf T)+\rho_ks_ks_k^\mathsf T,
\]
and the direction is $p_k=-H_kg_k$. L-BFGS stores only the last few
$(s_k,y_k)$ pairs and applies the inverse implicitly, making memory
$O(md)$ instead of $O(d^2)$.

\begin{theorem}[Local BFGS behavior]
For a twice continuously differentiable, strongly convex objective with
Lipschitz Hessian near $x^*$, BFGS with a suitable Wolfe line search converges
superlinearly under the standard Dennis--Moré conditions.
\end{theorem}

In computation, solve Newton systems rather than invert Hessians. Regularize
an indefinite Hessian using $H+\lambda I$, use a trust region, or switch to a
descent direction.

\section{Line Search}

Given a descent direction $p_k$ with $g_k^\mathsf Tp_k<0$, line search chooses
$\alpha_k>0$ and sets $x_{k+1}=x_k+\alpha_kp_k$. Exact line search minimizes
$\phi(\alpha)=f(x_k+\alpha p_k)$ but is usually unnecessary.

\begin{definition}[Armijo and Wolfe conditions]
The Armijo sufficient-decrease condition is
\[
 f(x_k+\alpha p_k)
 \leq f(x_k)+c_1\alpha g_k^\mathsf Tp_k,
 \qquad 0<c_1<1.
\]
The curvature condition is
\[
 \grad f(x_k+\alpha p_k)^\mathsf Tp_k
 \geq c_2g_k^\mathsf Tp_k,\qquad c_1<c_2<1.
\]
The strong Wolfe condition replaces the latter with an absolute-value bound.
\end{definition}

Backtracking starts from an initial step, often one, and repeatedly multiplies
it by $\rho\in(0,1)$ until Armijo holds.

\begin{algorithm}[H]
\caption{Backtracking line search}
\begin{algorithmic}[1]
\State Choose $\alpha\gets\alpha_0$, $\rho\in(0,1)$, $c_1\in(0,1)$.
\While{$f(x+\alpha p)>f(x)+c_1\alpha\grad f(x)^\mathsf Tp$}
  \State $\alpha\gets\rho\alpha$.
\EndWhile
\State Return $\alpha$.
\end{algorithmic}
\end{algorithm}

\begin{proposition}
If $\grad f$ is $L$-Lipschitz and $p$ is a descent direction, backtracking
terminates after finitely many reductions.
\end{proposition}

The Wolfe conditions prevent steps that are too short and support global
convergence theory for nonlinear conjugate-gradient and quasi-Newton methods.

\section{Subgradient Method}

For a convex, possibly nondifferentiable function $f$, a vector $g$ is a
subgradient at $x$ if
\[
    f(y)\geq f(x)+g^\mathsf T(y-x)\quad\text{for all }y.
\]
The set $\partial f(x)$ of all subgradients is the subdifferential.
For $f(x)=\abs x$, $\partial f(0)=[-1,1]$. A point minimizes a proper convex
function exactly when $0\in\partial f(x)$.

The subgradient iteration is
\[
    x_{k+1}=x_k-\alpha_kg_k,\qquad g_k\in\partial f(x_k).
\]
Unlike a gradient, $-g_k$ need not be a descent direction, so objective values
may oscillate.

\begin{theorem}[Subgradient bound]
Suppose $\norm{g_k}\leq G$ and $\norm{x_0-x^*}\leq R$. Then
\[
 \min_{0\leq j<k}\{f(x_j)-f(x^*)\}
 \leq\frac{R^2+G^2\sum_{j=0}^{k-1}\alpha_j^2}
           {2\sum_{j=0}^{k-1}\alpha_j}.
\]
Thus steps satisfying $\sum_k\alpha_k=\infty$ and
$\sum_k\alpha_k^2<\infty$ give convergence of the best objective value.
\end{theorem}

With a horizon known in advance, constant steps of order $1/\sqrt{k}$ give
the characteristic $O(k^{-1/2})$ rate. Proximal methods are usually much
faster when the nonsmooth term has a tractable proximal map.

\section{Proximal Gradient Method}

Consider the composite convex problem
\[
    \min_x F(x)=f(x)+g(x),
\]
where $f$ is differentiable and $L$-smooth while $g$ is convex and may be
nonsmooth.

\begin{definition}[Proximal operator]
For $\alpha>0$,
\[
 \operatorname{prox}_{\alpha g}(v)
 =\argmin_x\left\{g(x)+\frac1{2\alpha}\norm{x-v}^2\right\}.
\]
\end{definition}

The proximal-gradient step is
\begin{equation}
 x_{k+1}=\operatorname{prox}_{\alpha g}
             \{x_k-\alpha\grad f(x_k)\}.
 \label{eq:proximal-gradient}
\end{equation}
It minimizes a quadratic upper model of $f$ plus the exact $g$. Examples are
\[
 \operatorname{prox}_{\alpha\lambda\norm{\cdot}_1}(v)
 =\operatorname{sign}(v)(\abs v-\alpha\lambda)_+,
\qquad
 \operatorname{prox}_{\ind_C}(v)=\Pi_C(v).
\]

\begin{theorem}[Proximal-gradient rate]
If $F$ is convex, $f$ is $L$-smooth, and a minimizer exists, then with
$\alpha=1/L$,
\[
    F(x_k)-F(x^*)\leq\frac{L\norm{x_0-x^*}^2}{2k}.
\]
\end{theorem}

The gradient mapping
$G_\alpha(x)=\alpha^{-1}[x-\operatorname{prox}_{\alpha g}
(x-\alpha\grad f(x))]$ generalizes the gradient: $G_\alpha(x)=0$ exactly at
a first-order solution.

\section{Nesterov’s Accelerated Gradient Method}

Acceleration combines a gradient step with momentum. For convex $L$-smooth
$f$, one common form is
\[
 x_{k+1}=y_k-\frac1L\grad f(y_k),\qquad
 y_{k+1}=x_{k+1}+\frac{k}{k+3}(x_{k+1}-x_k).
\]

\begin{theorem}[Accelerated rate]
Nesterov's method can be chosen so that
\[
    f(x_k)-f(x^*)\leq\frac{2L\norm{x_0-x^*}^2}{(k+1)^2}.
\]
For $\mu$-strongly convex $f$, accelerated methods achieve a linear factor of
order $1-1/\sqrt{\kappa}$ rather than $1-1/\kappa$.
\end{theorem}

FISTA applies the same idea to proximal gradient and achieves $O(k^{-2})$ for
composite convex objectives. Momentum can cause nonmonotone objectives and
oscillation. Adaptive restart resets momentum when the objective or gradient
geometry indicates overshooting.

\section{Frank-Wolfe Method}

Frank--Wolfe solves $\min_{x\in C}f(x)$ for a compact convex set $C$ using a
linear minimization oracle rather than projection:
\[
 s_k\in\argmin_{s\in C}\grad f(x_k)^\mathsf Ts,\qquad
 x_{k+1}=(1-\gamma_k)x_k+\gamma_ks_k.
\]
It is useful when linear optimization over $C$ is cheap but Euclidean
projection is expensive, as for nuclear-norm balls.

The Frank--Wolfe gap
\[
    g_{\mathrm{FW}}(x)=
    \max_{s\in C}\grad f(x)^\mathsf T(x-s)
\]
is computable and satisfies $f(x)-f(x^*)\leq g_{\mathrm{FW}}(x)$ for convex
$f$. It is both a stopping criterion and a certificate.

\begin{theorem}[Frank--Wolfe rate]
If $f$ is convex and $L$-smooth on compact $C$ with diameter $D$, then with
$\gamma_k=2/(k+2)$,
\[
    f(x_k)-f(x^*)\leq\frac{2LD^2}{k+2}.
\]
\end{theorem}

Iterates are convex combinations of oracle outputs, often giving sparse or
low-rank solutions early. Standard Frank--Wolfe can zig-zag near a boundary;
away-step variants achieve faster rates under additional geometry.

\section{Proximal Point Algorithm}

The proximal point algorithm for a proper closed convex $f$ is
\begin{equation}
    x_{k+1}=\operatorname{prox}_{\lambda_k f}(x_k)
    =\argmin_x\left\{f(x)+\frac1{2\lambda_k}\norm{x-x_k}^2\right\}.
    \label{eq:ppa}
\end{equation}
Its optimality condition is
$0\in\partial f(x_{k+1})+\lambda_k^{-1}(x_{k+1}-x_k)$, or
$x_k-x_{k+1}\in\lambda_k\partial f(x_{k+1})$.

\begin{theorem}[Fejér monotonicity]
If $x^*$ minimizes $f$, then every exact proximal point step satisfies
\[
 \norm{x_{k+1}-x^*}^2
 \leq\norm{x_k-x^*}^2-\norm{x_{k+1}-x_k}^2.
\]
If $\inf_k\lambda_k>0$ and a minimizer exists, the iterates converge to a
minimizer in finite dimensions.
\end{theorem}

PPA is robust but each step may be as difficult as the original problem. Its
importance is conceptual: proximal gradient, augmented Lagrangian, and ADMM
can be interpreted as tractable approximations or operator-splitting forms of
proximal point iterations.

\section{Augmented Lagrangian Method}

For equality constraints
\[
    \min_x f(x)\quad\text{subject to}\quad Ax=b,
\]
the augmented Lagrangian is
\begin{equation}
 \mathcal L_\rho(x,y)
 =f(x)+y^\mathsf T(Ax-b)+\frac\rho2\norm{Ax-b}^2.
 \label{eq:augmented-lagrangian}
\end{equation}
The method of multipliers alternates
\[
 x_{k+1}\in\argmin_x\mathcal L_\rho(x,y_k),\qquad
 y_{k+1}=y_k+\rho(Ax_{k+1}-b).
\]
The quadratic penalty improves feasibility without requiring
$\rho\to\infty$, while the dual variable accumulates constraint violations.

Under convexity and standard constraint qualifications, a pair $(x^*,y^*)$
is optimal when
\[
    0\in\partial f(x^*)+A^\mathsf Ty^*,\qquad Ax^*=b.
\]
These KKT residuals should be monitored. Inner minimization need not be exact,
but its tolerances must become sufficiently accurate for convergence theory.

\section{Alternating Direction Method of Multipliers}

ADMM addresses
\[
    \min_{x,z}f(x)+g(z)\quad\text{subject to}\quad Ax+Bz=c.
\]
Using the scaled dual variable $u=y/\rho$, the iterations are
\begin{align*}
 x_{k+1}&=\argmin_x\left\{
 f(x)+\frac\rho2\norm{Ax+Bz_k-c+u_k}^2\right\},\\
 z_{k+1}&=\argmin_z\left\{
 g(z)+\frac\rho2\norm{Ax_{k+1}+Bz-c+u_k}^2\right\},\\
 u_{k+1}&=u_k+Ax_{k+1}+Bz_{k+1}-c.
\end{align*}
Separating $x$ and $z$ can turn a difficult joint problem into two simple
subproblems.

\begin{definition}[ADMM residuals]
The primal residual is
$r_{k+1}=Ax_{k+1}+Bz_{k+1}-c$. In the common constraint $x=z$, the dual
residual is $s_{k+1}=\rho(z_{k+1}-z_k)$, up to the appropriate matrix factor
in the general form.
\end{definition}

Stop only when both residuals are small relative to absolute and relative
tolerances. A large $\rho$ emphasizes feasibility but can slow dual progress;
a small $\rho$ does the reverse. Residual balancing adapts $\rho$ while
rescaling $u$ consistently.

\begin{theorem}[Convex ADMM convergence]
If $f$ and $g$ are proper, closed, convex, a saddle point exists, and the
subproblems are solved exactly, then the primal residual converges to zero,
objective values converge to the optimum, and dual variables converge under
standard rank conditions.
\end{theorem}

\section{Coordinate Descent Method}

Coordinate descent updates one block at a time:
\[
 x_j^{k+1}\in\argmin_t
 f(x_1^{k+1},\ldots,x_{j-1}^{k+1},t,x_{j+1}^k,\ldots,x_d^k).
\]
Updates may be cyclic, randomized, greedy, or performed in blocks. For a
smooth objective with coordinate Lipschitz constant $L_j$, a coordinate
gradient step is
$x_j\leftarrow x_j-L_j^{-1}\partial_jf(x)$.

\begin{theorem}[Randomized coordinate rate]
For convex smooth $f$ with coordinate-wise Lipschitz gradients, randomized
coordinate descent has an expected $O(1/k)$ objective gap. Under strong
convexity it converges linearly in expectation, with a rate controlled by the
coordinate smoothness constants.
\end{theorem}

Coordinate methods are effective when updates are closed form and data are
sparse. For lasso, each update is soft thresholding. Cache residuals rather
than recomputing $X\beta$ after each coordinate. Pure coordinate minimization
can fail for coupled nonsmooth nonconvex objectives, so assumptions matter.

\section{Stochastic Gradient Descent}

For finite-sum or expectation objectives,
\[
 f(x)=\frac1n\sum_{i=1}^nf_i(x)
 \quad\text{or}\quad
 f(x)=\E_\xi[F(x;\xi)],
\]
stochastic gradient descent (SGD) uses
\begin{equation}
    x_{k+1}=x_k-\alpha_kg_k,\qquad
    \E[g_k\mid x_k]=\grad f(x_k).
    \label{eq:sgd}
\end{equation}
A minibatch average lowers gradient variance roughly inversely with batch
size until correlations and hardware effects dominate.

\begin{theorem}[Strongly convex SGD rate]
Suppose $f$ is $\mu$-strongly convex, stochastic gradients are unbiased with
bounded second moment, and $\alpha_k$ is of order $1/k$. Then
\[
    \E[f(\bar x_k)-f(x^*)]=O(k^{-1})
\]
for a suitable averaged iterate. For convex but not strongly convex objectives,
the standard rate is $O(k^{-1/2})$.
\end{theorem}

A constant learning rate reaches a noise neighborhood rather than converging
exactly, but is useful for rapid initial progress. Decay reduces the
steady-state variance. Momentum updates
\[
    v_{k+1}=\beta v_k+g_k,\qquad x_{k+1}=x_k-\alpha v_{k+1}
\]
smooth noisy directions and accelerate persistent ones.

Shuffle finite datasets between epochs. Report gradient evaluations or data
passes, not iterations alone, when comparing batch sizes. For nonconvex
smooth objectives, a common result controls the average squared gradient norm,
so convergence means approaching stationarity rather than a global minimum.

\section{Adaptive Methods}

Adaptive methods rescale coordinates using past gradients:
\[
    x_{k+1}=x_k-\alpha D_k^{-1}m_k.
\]
They are helpful for sparse gradients and poorly scaled coordinates.
However, their implicit geometry, memory cost, and generalization can differ
from SGD. The stabilizer $\varepsilon$ is part of the algorithm, not merely a
guard against division by zero.

\subsection{ADAM}

Adam maintains exponential first and second moments:
\begin{align*}
 m_k&=\beta_1m_{k-1}+(1-\beta_1)g_k,\\
 v_k&=\beta_2v_{k-1}+(1-\beta_2)g_k^{\odot2},\\
 \widehat m_k&=m_k/(1-\beta_1^k),\qquad
 \widehat v_k=v_k/(1-\beta_2^k),\\
 x_{k+1}&=x_k-\alpha\frac{\widehat m_k}
                    {\sqrt{\widehat v_k}+\varepsilon}.
\end{align*}
Bias correction matters early because moments start at zero.

Vanilla Adam does not converge for every convex online problem; AMSGrad
repairs a problematic decrease in the effective second-moment accumulator by
using a coordinatewise running maximum. In deep learning, AdamW decouples
weight decay from the adaptive gradient update. Adding $\lambda x$ to the
gradient is not equivalent to multiplicative weight decay under a
coordinate-dependent preconditioner.

\subsection{AdaGrad}

Diagonal AdaGrad uses
\[
    v_k=v_{k-1}+g_k^{\odot2},\qquad
    x_{k+1}=x_k-\alpha\frac{g_k}{\sqrt{v_k}+\varepsilon}.
\]
Coordinates with frequent large gradients receive smaller future steps.

\begin{theorem}[AdaGrad regret, informal]
For convex online losses and a bounded domain, diagonal AdaGrad has regret
controlled by
\[
    \sum_{j=1}^d
    \sqrt{\sum_{k=1}^Tg_{k,j}^2},
\]
up to a factor depending on the domain diameter. This can be much smaller
than a bound based on the worst gradient when gradients are sparse.
\end{theorem}

Because $v_k$ only increases, learning rates may decay too aggressively in a
long nonstationary run. RMSProp and Adam use exponential rather than cumulative
second moments.

\section{Shampoo and Muon}

Diagonal adaptive methods ignore correlation between parameter coordinates.
Full-matrix AdaGrad captures it but requires prohibitive storage and matrix
roots. Shampoo exploits tensor structure. For a matrix gradient
$G_k\in\R^{m\times n}$, a basic form accumulates
\[
    L_k=\varepsilon I_m+\sum_{t\leq k}G_tG_t^\mathsf T,\qquad
    R_k=\varepsilon I_n+\sum_{t\leq k}G_t^\mathsf TG_t
\]
and preconditions by
\begin{equation}
    \widetilde G_k=L_k^{-1/4}G_kR_k^{-1/4}.
    \label{eq:shampoo}
\end{equation}
Higher-order tensors receive one preconditioner per mode. Matrix roots are
expensive, so practical implementations update them intermittently, use
blocking, and combine them with grafting or momentum. Damping controls nearly
singular eigendirections.

Muon is designed primarily for matrix-valued hidden-layer weights. It forms a
momentum matrix $M_k$ and approximately replaces its singular values by one,
\[
    M_k=U\Sigma V^\mathsf T
    \quad\longmapsto\quad
    \operatorname{polar}(M_k)=UV^\mathsf T,
\]
usually through a small number of Newton--Schulz matrix iterations rather than
an SVD. The resulting direction is then scaled and used for the parameter
update. Vector parameters and some embedding or output parameters are commonly
handled by another optimizer.

Shampoo uses accumulated second-moment geometry; Muon orthogonalizes a current
momentum direction. They are therefore not interchangeable despite both using
matrix operations. Exact scaling conventions vary by implementation, so a
reproducible description must specify normalization, Newton--Schulz polynomial,
iteration count, momentum, weight decay, and parameter groups.

\section{Automatic Differentiation}

Automatic differentiation (AD) applies the chain rule to a computational
graph. It is neither symbolic differentiation nor finite differences.
For $f:\R^n\to\R^m$ with Jacobian $J_f(x)$:
\begin{itemize}
 \item forward mode computes a Jacobian--vector product $J_f(x)v$;
 \item reverse mode computes a vector--Jacobian product
 $u^\mathsf TJ_f(x)$.
\end{itemize}
Forward mode is attractive when inputs are few; reverse mode is attractive
for a scalar loss with many parameters, because one reverse sweep obtains the
full gradient.

\begin{example}[Reverse-mode chain rule]
If $z=Wx$, $a=\tanh z$, and $\ell=\tfrac12\norm{a-y}^2$, reverse mode computes
\[
 \bar a=a-y,\qquad
 \bar z=\bar a\odot(1-\tanh^2z),\qquad
 \grad_W\ell=\bar z x^\mathsf T.
\]
The barred variables are adjoints: derivatives of the final loss with respect
to intermediate values.
\end{example}

Reverse mode stores intermediate activations, trading memory for computation.
Checkpointing recomputes selected activations during the backward pass.
Higher-order derivatives are obtained by composing modes; Hessian--vector
products can be computed without forming the Hessian.

Common bugs arise from in-place mutation, unintended detachment of the graph,
nondifferentiable control flow, and unstable primitives. Gradient checking
with centered finite differences is useful on small problems, but step sizes
that are too small suffer cancellation and those too large suffer truncation
error.

\section{Applications}

Optimization algorithms should be chosen from problem structure: smoothness,
convexity, separability, sparsity, constraint geometry, and data-access cost.
The following examples connect the general methods to recurring machine
learning and inverse-problem formulations.

\subsection{Matrix Factorization}

Given a partially observed matrix $M$, low-rank factorization solves
\begin{equation}
 \min_{U\in\R^{m\times r},\,V\in\R^{n\times r}}
 \frac12\sum_{(i,j)\in\Omega}
 (M_{ij}-u_i^\mathsf Tv_j)^2
 +\frac\lambda2(\norm U_F^2+\norm V_F^2).
 \label{eq:matrix-factorization}
\end{equation}
The problem is nonconvex jointly but convex in either factor with the other
fixed. Alternating least squares exploits this biconvex structure. SGD samples
an observed entry and updates only $u_i$ and $v_j$:
\[
 e_{ij}=M_{ij}-u_i^\mathsf Tv_j,\quad
 u_i\leftarrow u_i+\alpha(e_{ij}v_j-\lambda u_i),\quad
 v_j\leftarrow v_j+\alpha(e_{ij}u_i-\lambda v_j),
\]
where the second update should use either the old $u_i$ or a documented
sequential convention.

The factors are nonidentifiable: $UV^\mathsf T=(UA)(VA^{-\mathsf T})^\mathsf T$
for invertible $A$. Prediction of the product may be well defined even when
individual factors are not. Random initialization, regularization, and
validation of rank $r$ are essential.

The convex alternative minimizes a data-fit term plus nuclear norm
$\norm X_*=\sum_j\sigma_j(X)$. Its proximal operator performs singular-value
soft thresholding, while Frank--Wolfe on a nuclear-norm ball adds rank-one
directions using leading singular vectors.

\subsection{Support Vector Machines}

The linear soft-margin SVM has a smooth quadratic regularizer and nonsmooth
hinge losses:
\[
 \min_{\beta_0,\beta}
 \frac\lambda2\norm\beta^2+
 \frac1n\sum_{i=1}^n[1-y_i(\beta_0+x_i^\mathsf T\beta)]_+.
\]
Subgradient methods are simple but slow. Coordinate descent is effective in
the dual; stochastic subgradient descent scales to large sparse primal
problems. For one observation, a subgradient with respect to $\beta$ is
\[
 \lambda\beta-
 y_ix_i\ind\{y_i(\beta_0+x_i^\mathsf T\beta)<1\},
\]
with a set-valued choice exactly at margin one.

Squared hinge loss is differentiable once but changes the robustness and dual.
Kernel methods favor dual solvers, while high-dimensional sparse linear SVMs
usually favor primal methods. Optimization tolerance should be judged relative
to statistical error: solving the training objective far beyond the precision
supported by the data wastes computation.

\subsection{Compressed Sensing}

Compressed sensing recovers a sparse vector $x_0\in\R^d$ from
$b=Ax_0+\varepsilon$ with fewer measurements than ambient coordinates.
The ideal problem
\[
    \min_x\norm{x}_0\quad\text{subject to}\quad Ax=b
\]
is combinatorial. Basis pursuit replaces cardinality by the convex $\ell_1$
norm:
\begin{equation}
    \min_x\norm x_1\quad\text{subject to}\quad Ax=b.
    \label{eq:basis-pursuit}
\end{equation}
With noise, common forms are
\[
 \min_x\frac12\norm{Ax-b}^2+\lambda\norm x_1
 \quad\text{or}\quad
 \min_x\norm x_1\ \text{subject to }\norm{Ax-b}\leq\delta.
\]
ISTA, FISTA, ADMM, and coordinate descent apply directly.

\begin{definition}[Restricted isometry property]
A matrix $A$ has restricted isometry constant $\delta_s$ if $\delta_s$ is the
smallest number such that
\[
 (1-\delta_s)\norm x^2\leq\norm{Ax}^2
 \leq(1+\delta_s)\norm x^2
\]
for every $s$-sparse $x$.
\end{definition}

\begin{theorem}[Sparse recovery, representative form]
If $A$ satisfies a sufficiently strong restricted isometry condition, such as
$\delta_{2s}<\sqrt2-1$, then basis pursuit exactly recovers every $s$-sparse
vector from noiseless measurements. With noise and approximate sparsity, the
reconstruction error is bounded by a noise term plus the best $s$-term
approximation error.
\end{theorem}

Random sub-Gaussian measurement matrices satisfy suitable RIP conditions with
roughly $s\log(d/s)$ measurements. Feature standardization and the choice of
$\lambda$ remain important in practical statistical problems.

\subsection{Deep Learning}

A neural network is a composition
\[
 h_0=x,\qquad
 z_\ell=W_\ell h_{\ell-1}+b_\ell,\qquad
 h_\ell=\sigma_\ell(z_\ell),
\]
trained by minimizing an empirical loss, usually with regularization. Reverse-
mode AD is backpropagation applied to this graph.

The objective is nonconvex and has symmetries: hidden units may be permuted,
and homogeneous networks allow compensating rescalings across layers.
Consequently parameter identifiability and Hessian positive definiteness are
not expected. Practical optimization seeks a useful low-loss stationary
region, not a certified global minimum.

Initialization controls signal and gradient propagation. Xavier scaling is
suited to approximately symmetric activations; He scaling is suited to ReLU-
like activations. Normalization, residual connections, and careful parameter
scaling improve conditioning. Gradient clipping controls rare large updates
but changes the optimization direction.

Minibatch SGD and Adam-type methods dominate because full gradients are
expensive. A typical training loop separates:
\begin{enumerate}
 \item forward evaluation and loss computation;
 \item reverse-mode gradient computation;
 \item optional clipping or gradient accumulation;
 \item optimizer and learning-rate-schedule updates;
 \item validation under evaluation-mode behavior.
\end{enumerate}

Weight decay, data augmentation, dropout, and early stopping are
regularization choices as well as optimization choices. Training loss alone
cannot choose them. Reproducible computation records random seeds, data order,
precision, optimizer conventions, schedules, and checkpoint selection.

Mixed-precision training increases throughput but can underflow small
gradients; loss scaling and higher-precision accumulators address this.
Distributed data-parallel training averages gradients across workers. Very
large batches reduce gradient noise but often require learning-rate and
warmup adjustments.

\begin{remark}
For both classical and deep models, optimization error, estimation error, and
approximation error are distinct. A lower training objective does not imply a
lower population risk, and an optimizer comparison is meaningful only when
data processing, model, tuning budget, and stopping rules are controlled.
\end{remark}


\end{document}
